% Électronégativités (échelle de Pauling) des éléments des trois premières lignes (L'essentiel, chapitre 6)
\begin{tikzpicture}[x=1cm,y=1cm]
\foreach \col/\lig/\sym/\ch in {1/1/H/2.20,1/2/Li/0.98,2/2/Be/1.57,3/2/B/2.04,4/2/C/2.55,5/2/N/3.04,6/2/O/3.44,7/2/F/3.98,
  1/3/Na/0.93,2/3/Mg/1.31,3/3/Al/1.61,4/3/Si/1.90,5/3/P/2.19,6/3/S/2.58,7/3/Cl/3.16}{
  \pgfmathsetmacro\teinte{round(8+18*(\ch-0.9))}
  \pgfmathsetmacro\clair{\ch<2.9 ? 1 : 0}
  \ifdim\clair pt>0.5pt\def\coul{otInk}\else\def\coul{white}\fi
  \filldraw[fill=otNavy!\teinte,draw=white,line width=1pt,rounded corners=1mm] ({1.55*(\col-1)},{-1.25*(\lig-1)}) rectangle ++(1.45,-1.15);
  \node[text=\coul,font=\otSemi\fontsize{10}{11}\selectfont] at ({1.55*(\col-1)+0.725},{-1.25*(\lig-1)-0.42}) {\sym};
  \node[text=\coul,font=\otLight\fontsize{8}{9}\selectfont] at ({1.55*(\col-1)+0.725},{-1.25*(\lig-1)-0.83}) {\num{\ch}};}
\draw[otGreenDark,line width=1.1pt,-{Stealth[length=2.6mm]}] (1.7,0.25) -- (10.75,0.25)
  node[midway,above,pcLegende,text=otGreenDark] {l'électronégativité augmente de gauche à droite};
\draw[otGreenDark,line width=1.1pt,-{Stealth[length=2.6mm]}] (11.2,-3.75) -- (11.2,-1.3)
  node[midway,right,pcLegende,text=otGreenDark,align=left] {et de bas\\en haut};
\end{tikzpicture}
\par\smallskip\legende{Les gaz nobles, qui ne forment pas de liaisons dans les conditions usuelles, n'ont pas d'électronégativité dans cette échelle.}
